% ============================================================
%  Poul Martin Møller,
%  "Forelæsnings-Paragrapher over Moralphilosophien"
%
%  Numbered lecture-paragraphs (§§) on moral philosophy. Per the
%  editor's title-note: "Efter Collegier fra 1837, sammenlignede
%  med Forfatterens Manuscript. F. C. O." (F.C. Olsen) — i.e.
%  reconstructed from lecture notes of 1837 collated with the
%  author's manuscript.
%
%  Published posthumously in:
%  Efterladte Skrifter af Poul M. Møller, bd. 5
%  (Kjøbenhavn: C.A. Reitzel, 1855--56), pp. 141--164.
%  Transcribed from that edition's Fraktur setting.
%  PDF (bibliotek scan, efterladte-skrifter5.pdf) pp. 152--175;
%  offset PDF = printed + 11.
%
%  TRANSCRIPTION (Danish)
%  Conventions: original orthography kept (aa-spellings,
%  capitalised nouns, æ/ø, "Characteer", "Conseqvents",
%  the -mus endings Determinismus/Indifferentismus as printed).
%  Fraktur long-s -> s. Letterspaced emphasis -> \emph{}.
%  The editor's "Anm." remarks are set smaller. Printed page
%  numbers -> \opage{N}.
% ============================================================

\documentclass[12pt]{article}
\usepackage[utf8]{inputenc}
\usepackage[T1]{fontenc}
\usepackage{graphicx}
\DeclareUnicodeCharacter{0254}{\reflectbox{c}}
\usepackage{libertinus}
\usepackage{libertinust1math}
\usepackage{textalpha}
\usepackage[danish]{babel}
\usepackage[protrusion=true,expansion=true]{microtype}
\usepackage[margin=1.4in]{geometry}
\usepackage{setspace}
\usepackage{fancyhdr}
\usepackage{marginnote}
\usepackage[colorlinks=true,linkcolor=black,urlcolor=black,
  pdftitle={Forelæsnings-Paragrapher over Moralphilosophien},
  pdfauthor={Poul Martin Møller}]{hyperref}
\setlength{\headheight}{15pt}
\pagestyle{fancy}
\fancyhf{}
\fancyhead[LE,RO]{\thepage}
\fancyhead[RE]{\textit{Moralphilosophien}}
\fancyhead[LO]{\textit{Poul Martin Møller}}
\setlength{\parindent}{1.5em}
\setlength{\parskip}{0pt}
\onehalfspacing

\newcommand{\opage}[1]{\marginnote{\footnotesize\textit{[#1]}}}
\newcommand{\ombreak}{\par\medskip\begin{center}\rule{0.32\textwidth}{0.4pt}\end{center}\medskip}
\newcommand{\slaw}[1]{\par\medskip\begin{center}\textbf{\S~#1.}\end{center}\nopagebreak}
\newcommand{\anm}[1]{\par\smallskip{\small\hangindent=1.6em\hangafter=0\textit{Anm.}\quad #1\par}\smallskip}
\newcommand{\mchap}[2]{\par\bigskip\begin{center}\large\textbf{#1}\\[0.3em]\textbf{#2}\end{center}\medskip}

\title{\textbf{Forelæsnings-Paragrapher\\[0.2em] over Moralphilosophien}\footnote{Efter Collegier fra 1837, sammenlignede med Forfatterens Manuscript. F.\ C.\ O.}}
\author{Poul Martin Møller}
\date{Efterladte Skrifter, bd.\ 5}

\begin{document}
\maketitle

\slaw{1}
\opage{141}Efter den almindelige Forestillingsmaade er Moralen Videnskaben om Menneskets fornuftige Handlinger og det Øiemed, der ved dem skal opnaaes. Den Erkjendelse af det Gode, som i Moralphilosophien bliver videnskabelig udviklet, er først tilstede paa en umiddelbar Maade, saa at den Maxime, der i de enkelte Tilfælde bør følges, bliver indseet uden Bevidsthed om nogen Grund derfor. Hvor Spørgsmaalet om de practiske Reglers Grund indtræder, er det en naturlig Tankegang at søge det sidste Øiemed for Menneskenes Bestræbelser og deraf at udlede alle de specielle Regler for dem. Efter den dogmatiske Methode, som fremkommer derved, udtrykkes den øverste practiske Regel i een Formel, der betragtes som indlysende ved sig selv og lægges til Grund for det hele System. Paa denne Maade have flere modsatte Principer under Videnskabens historiske Udvikling været gjorte gjeldende mod hinanden, og de vigtigste Pligter og Dyder have med større eller mindre Eensidighed været udviklede af ethvert af dem, saa at intet af dem kan være aldeles urigtigt. Men ved at underkaste dem en videnskabelig Critik bragte man deres Eensidighed til Bevidsthed og kom til det Resultat, at Menneskets fornuftige Stræben ei lader sig bestemme ved nogen enkelt almeengyldig Regel. I den egentlig philosophiske Methode ere disse to Synsmaader forenede med hinanden, idet Mangelen i ethvert eensidigt Princip bliver efterviist, men derved tillige gjøres Overgang til en ny Bestemmelse for Principet, saa at Videnskabens sande Princip ud\opage{142}vikles ved den Tænkning, der gjennemgaaer alle de eensidige Bestemmelser og følger deres naturlige Overgang i hinanden.

\slaw{2}
I Moralen forudsættes, at Mennesket er et frit Væsen, men Frihedens Begreb erholder nye Bestemmelser i Videnskaben selv. --- Den Frihed, som forudsættes, er ei den, der antages af Indifferentisterne, der paastaae, at Mennesket i ethvert Moment kan bestemme sig imod alle Motiver. Disse forsvare deres Theorie deels med en foregiven umiddelbar Bevidsthed om, at man undertiden virkelig bestemmer sig mod alle Motiver; deels ved den Paastand, at Mennesket, naar Motiverne ere lige stærke, ei kan bestemme sig uden med absolut Vilkaarlighed; de mene endelig, at den enkelte Handling kun efter deres Synsmaade kan have moralsk Værd eller Uværd. Men herpaa maa svares, paa det Første, at umiddelbar Erfaring ikke kan lære noget Negativt; paa det Andet, at den forudsatte absolute Ligevægt ei kan finde Sted; paa det Tredie, at Handlingen kun ved en urigtig Abstraction kan bedømmes som adskilt fra Characteren. Desuden er Indifferentismen, forfulgt i sine Conseqventser, aldeles utænkelig og strider, idet den ophæver Begrebet Character, mod Erfaringens udtrykkelige Vidnesbyrd.

\slaw{3}
Paa den anden Side er den eensidige Determinismus ogsaa urigtig, naar den enten gjør Villien aldeles afhængig af Erkjendelsens videnskabelige Udvikling, eller naar den antager, at den enkelte Handling er et nødvendigt Product af Characteren og de udvortes Omstændigheder, og Characteren et Resultat af alle de foregaaende Handlinger. Hvis den da anseer den første Handling for fri (Prædeterminismus), da er den heri inconseqvent; thi de væsentlige Vanskeligheder, som finde Sted ved at begribe de følgende Handlingers Frihed, gjelde ligesaa fuldt ved den første Handling. Antages derimod ogsaa den første Handling for ufri, \opage{143}er Determinismen ikke væsentlig forskjellig fra Fatalismen, der strider saavel mod den moralske Bevidsthed, som mod det ægte Begreb om en skabende Guddom.

\slaw{4}
Efter det rigtige Begreb om Frihed tænkes den enkelte Handling ei at være ganske uafhængig af \emph{den empiriske Characteer}, hvilket vilde hæve Begrebet om Dyd og Udyd; men heller ei at være aldeles bestemt derved; tvertimod kan den empiriske Characteer i moralsk Henseende forandres, dog saaledes, at Villiens forandrede Retning kun efterhaanden medfører Færdigheden i det Gode eller i det Onde. En usædelig Characteer kan aldrig være aldeles uforanderlig, men derimod kan en dydig Characteer være saa befæstet, at ingen Forandring af Grundsætninger kan foregaae deri. --- Motiver bestemme Handlingerne, men de ere aldrig blot udvortes. Kun fra den ene Side betragtede kunne de være blot udvortes, fra den anden ere de indvortes og komme til at gjelde ved Villiens bestemte Retning. Omvendelse kan vel gaae pludselig for sig, forsaavidt som forandrede Maximer med Eet kunne følges; men Sindelagets totale Forandring kommer kun i Stand ved megen Øvelse.

\slaw{5}
Moralsk Betydning have kun de Livsyttringer, som Mennesket har vidst af, og som det har villet. De forskjellige Grader af Frihed, som her kunne tænkes at finde Sted, udvikles i Moralens første Hoveddeel, \emph{Læren om Tilregnelse}, hvori Begrebet om en Handling bestemmes, uden at der endnu tages Hensyn til, hvad den gaaer ud paa; i anden Hoveddeel, \emph{Læren om Lyksalighed}, udvikles den menneskelige Stræbens forskjellige Formaal, og, da disse komme til at stride mod hinanden, men fornuftigviis maae underordnes et almindeligt Øiemed, udkræves endnu en tredie Hoveddeel af Moralen, \emph{Læren om det sande Gode}.

\mchap{Første Capitel.}{Om Tilregnelse.}

\slaw{6}
\opage{144}Mennesket siges at foretage en Handling, naar det er frivillig Aarsag til en Forandring i Omverdenen. I Handlingens ufravigelige Bestemmelse, Frivilligheden, ere to Bestemmelser at adskille, 1) \emph{Selvbestemmelsen under den yttrede Virksomhed}, 2) \emph{Opfattelsen af de Objecter}, hvorpaa den gaaer ud; og disse to Bestemmelser falde i den blot frivillige Handling sammen i eet Moment. Mangler een af disse Bestemmelser, saa at Mennesket bevæges af en uimodstaaelig udvortes Magt, eller har urigtige Forestillinger om de Gjenstande, det virker paa, da ere de Virkninger, hvortil Mennesket er Aarsag, at ansee for ufrivillige. En Forbindelse af det Frivillige og Ufrivillige finder Sted i de blandede Gjerninger, som ere halv frivillige og halv ufrivillige. Disse finde Sted, \emph{deels} naar Mennesket tvungent udøver en Handling, som det ikke vil i og for sig selv, men kun for at undgaae et efter dets Mening større Onde, \emph{deels} naar det befinder sig i en Naturtilstand, som medfører en fordunklet Bevidsthed. --- Forsaavidt som alle Objecter for Handlingen for sig have en Beskaffenhed, der virker med til at bestemme Handlingens Beskaffenhed, er egentligen enhver menneskelig Handling blandet.

\anm{Til de Vildfarelser, der gjøre Menneskenes Handlinger ufrivillige, regnes ei Uvidenheden om Modsætningen mellem Godt og Ondt i Almindelighed.}

\slaw{7}
\opage{145}Naar Mennesket gjør den Erfaring, at de udvortes Omstændigheder kunne opfattes urigtigt og hindre Villiens Udførelse, udvikles dets Bevidsthed saaledes, at det først fatter et Forsæt, hvori Handlingen har en indre Tilværelse i dets Bevidsthed, og først derefter udføre det i Sandseverdenen. I den forsætlige Handling lader sig adskille de tre Sider: \emph{Overlæg, Beslutning og Udførelse}. Enhver forsætlig Handling er frivillig, men en Handling kan være blot frivillig, uden at være forsætlig: naar den er forsætlig, er den mere tilregnelig, end som blot frivillig.

\slaw{8}
Men den Harmonie, som vi have tænkt os at finde Sted mellem Forsættet og Handlingens Udførelse, kan blive forstyrret derved, at de opfattede Omstændigheder kunne hænge sammen med andre ikke opfattede, og som enten forandre Handlingen saaledes, at den ikke svarer til Forsættet, eller give den uforudsete Følger. Naar Handlingens Forandring eller de uforudsete Følger ei kunde eller burde have været forudsete, finder et \emph{Tilfælde} Sted, som slet ikke kan tilregnes.

\slaw{9}
Naar Mennesket handler uden at indsee Muligheden af, at det ved sin Handling kan blive Skyld i en Begivenhed, skjøndt det kunde og burde forudsee det, da bliver det Uforsætlige tilregnet det, men dog kun som \emph{Forseelse}. Denne har sin Grund i en Forsømmelse af Selvvirksomhed, der drager en Uvidenhed eller Naturtilstand efter sig, hvorved den enkelte Feil i Bevidstheden, som altid er forbunden med Forseelsen, bliver foraarsaget. Nærmest er da Uvidenheden eller Naturtilstanden tilregnelig, og middelbart den Handling, som er foranlediget derved.

\slaw{10}
Mere tilregnelig end den blot forsætlige Handling er den \emph{overveiede}, hvortil man i Almindelighed kan henregne den efter Plan udførte. Denne finder tydelig Sted, naar en Handling ikke blot er forberedt ved flere foregaaende, ved hvilke den Handlende har sat sig i Stand til at udføre den, men Rækken af disse Handlinger er bleven afbrudt ved mellemkommende Handlinger i andre Hensigter.

\slaw{11}
\opage{146}Til Handling i Ordets vidtløftigste Betydning regner man ogsaa de aandelige Livsyttringer, Forestillinger og Tanker. Disse kunne undertiden ufrivilligen opstaae hos Mennesket efter Ideeassociationens Love, men selv de meest ufrivillige Handlinger ere dog afhængige af Menneskets Characteer. Endnu mere beroer det paa Menneskets Frihed, om det vil afbryde eller fortsætte den Forestillingsrække, som tilfældigviis paatrænges det. Den frivillige Reproduction af Forestillinger, Phantaseringen og Tænkningen kan og, ligesom Handlingen i engere Betydning eller den udvortes Handling, foretages saavel efter Forsæt som efter Plan.

\slaw{12}
Ved enhver forsætlig Handling har Mennesket altid en \emph{sidste Hensigt} eller et \emph{Formaal for Øie}: det udøver aldrig en Handling blot for at udøve den. Den nærmeste Hensigt ved en Handling kan vel være en anden Handling, som det ved den første vil gjøre mulig; ved den anden Handling kan det have en tredie til Hensigt o.\ s.\ v.; men Rækken af Handlinger gaaer dog altid tilsidst ud paa et Formaal. Disse Formaal udgjøre Handlingernes blot indre Side, og med disses Undersøgelse træde vi ind paa den egentlige Morals Enemærker.

\anm{Bevidstheden om, at Menneskets Handlinger altid ere Midler, har givet Anledning til den feilagtige Synsmaade, at Menneskets Villie skulde bedømmes blot efter dets Handlingers Formaal. Ingen Handling kan betragtes \emph{blot} som Middel: Summen af alle Menneskets Handlinger er det objective Billede af dets Villie.}

\mchap{Andet Capitel.}{Om Lyksalighed.}

\slaw{13}
Om Indholdet af det rette Formaal vide vi for det Første ikke Andet, end at der maa være et Formaal for Menneskets Be\opage{147}stræbelser, ɔ: at Mennesket maa søge personlig Tilfredsstillelse ved sine Handlinger eller have \emph{Interesse} for dem. Herved lades det endnu aldeles ubestemt, af hvad Art Interessen er, (om den f.\ Ex.\ gaaer paa den sandselige Begjerings Tilfredsstillelse eller paa Opnaaelsen af et fornuftigt Øiemed). Dette Begreb om Interesse, betragtet som det øverste Moralprincip, er det allermeest indholdsløse Princip, men udhæver dog en nødvendig Bestemmelse ved det Gode.\footnote{Forf.\ Msc.: ved Menneskets Øiemed. F.\ C.\ O.}

\anm{Forsaavidt som Tilfredsstillelsen, betragtet som Handlingernes Øiemed, forudsætter, at eet Savn efter det andet maa udfyldes, kan \emph{Virksomheden}, som det Blivende i disse Overgange, ogsaa ved en eensidig Abstraction betragtes som det adæqvate Udtryk for Menneskets Bestemmelse.}

\slaw{14}
Derved, at Mennesket i sine Handlinger udfører sin Interesse eller sit Livs Øiemed, tilveiebringes en Harmonie mellem dets Villie og Objecternes Verden, og den meest umiddelbare Maade, hvorpaa denne Overeensstemmelse kommer til den Handlendes Bevidsthed, er ved \emph{Følelsen}. Den \emph{behagelige Følelse} bringes man derfor ved en naturlig Tankegang til at ansee for den menneskelige Virksomheds Formaal, og denne bliver nu Gjenstanden for vor Betragtning. Subjectets Harmonie med Omverdenen tænkes som Noget, der skal tilveiebringes og altsaa endnu ikke finder Sted. Mennesket svarer i sin Stilling til Omverdenen ei til sin Bestemmelse: det har den for Øie, som et objectivt Formaal; men dets Overeensstemmelse eller Uovereensstemmelse hermed kjende vi endnu kun som det Behagelige eller Ubehagelige. Principet bliver altsaa, at skye det Ubehagelige, eller hvad der er det Samme, at stræbe efter det Behagelige.

\slaw{15}
Imellem de Moralister, der sætte Livets ægte Lyst i Tilfreds\opage{148}stillelsen af de naturlige Begjeringer, optræder den \emph{kyrenaiske Skole} med det Princip, at concentrere den størst mulige Lyst i hvert enkelt Moment af Livet og at sætte Menneskets Bestemmelse i en uafbrudt Række af enkelte Fornøielser. Men denne Regel ophæver sig selv ved den mindste Eftertanke, da denne fører til at tage Lystens Følger i Betragtning og undertiden at fornegte sig en Lyst, for at undgaae en Smerte eller for at opnaae en større Lyst, ɔ: at overtræde Principet for Principets Skyld. Men har man først indrømmet \emph{Tanken} sin Ret og indseet, at den leder Mennesket til, af Klogskab undertiden at underkaste sig en Ulyst, da erkjendes let, at det ei staaer i de Dødeliges Magt, at forskaffe sig en overveiende Sum af behagelige Momenter, og det er conseqvent, derfra at gaae over til det \emph{kyniske Princip}: at sikkre sig Lyksalighed ved en Uafhængighed, der opnaaes ved Tilfredsstillelsen af saa faa Begjeringer, som muligt. Disse to modsatte Retninger forenedes af Epikur, idet han vel erklærede Lysten for at være det høieste Formaal, men dog deels indsaae Selvfornegtelsens Nødvendighed, deels begyndte at erkjende den aandelige Lysts Betydning. Den \emph{aandelige Lyst}, som er den sande menneskelige Lyst, kan ikke sættes i den øieblikkelige Tilfredsstillelse af de hinanden modstridende Drifter, men kun i Udførelsen af eet almindeligt Formaal, hvorved der bringes Eenhed i den menneskelige Virksomhed. Da den Virksomhed, der fremmer det rette Øiemed, maa være uadskillelig forbunden med en Lyst, indeholder og det \emph{eudaimonistiske Princip}, rigtig forstaaet, een af Moralprincipets vigtigste Bestemmelser. Men til at finde den almindelige Retning af Villien, der medfører Lyst, have vi endnu ingen anden Veiledning end \emph{Følelsen}, der paa dette Trin viser sig som \emph{moralsk Følelse}.

\slaw{16}
Naar den \emph{moralske Følelse} antages for Handlingernes rette Princip, er det ikke den momentane Tilfredsstillelse af de \opage{149}sandselige Begjeringer, der betragtes som Lystens Kilde, men en almindelig Retning af Villien antages for at være den rigtige, fordi den i og for sig behager Individet, enten det naaer sit Formaal eller ikke. I dette Princip ligger den Sandhed, at den ægte moralske Villie saaledes maa være et Udtryk af Personlighedens inderste Væsen, at Mennesket befinder sig vel ved dens Fordringer. Deri ligger, at sand Moralitet er uadskillelig forbunden med en særegen Lyst. Den sande Moralitet maa følgelig ogsaa have Følelsens Form, men denne Følelses Indhold er endnu ikke bekjendt. Denne sidste Omstændighed har foranlediget Andre til at ansee Indholdet for at være givet aldeles udenfra, ved \emph{Opdragelsen og Statens Indretninger og Vedtægter}. Men ved denne Mening, som egentlig blot sætter Problemets Opløsning længer ud, tilintetgjøres aldeles den moralske Frihed, ifølge hvilken Mennesket selv kan bestemme sit Livs Øiemed. --- Spørges der nærmere om det Blivende og Almindelige i den Villiens Retning, som behager Mennesket, da kan det for det Første sikkert betragtes som en Bestræbelse for at bevare dette egne Jeg, og Mange have derfor antaget \emph{Selvbevaringen} for Livets Øiemed.

\slaw{17}
\emph{Selvbevaringen} som Principet for Menneskets Handlinger bringer det til at søge Lysten som Middel for Selvbevaringen. Det betragtes paa dette Standpunct som Livets Øiemed, at tilveiebringe og benytte et saadant Forhold mellem Subjectet og Omverdenen, at alt Udvortes tjener til det individuelle Livs Fremme. I sin laveste Skikkelse er Selvbevaringen den \emph{physiske Selvopholdelse} og viser sig deri, at Mennesket stræber at sikkre sin legemlige Tilværelse mod Angreb udenfra og skaffe sig Næring og Incitament for det blot organiske Livs Fortsættelse. Dernæst viser den sig som en Stræben efter \emph{Besiddelse} af de Naturting, som ere Betingelser for dets physiske Existents, eller Stræben efter \opage{150}Eiendom. Ved Stræben efter Eiendom kommer Mennesket til ikke blot at staae i Forhold til Sandseverdenen, men til andre Mennesker, som erkjende dets Ret til, hvad det besidder. Den Anerkjendelse af Subjectets Villie, der finder Sted i Eiendom, er Begyndelsen til dets udvortes \emph{Ære} eller dets Agtelse hos andre Mennesker, der ogsaa for sig kan gjøres til Formaal. (En conseqvent Fortsættelse af Stræben efter Ære er saavel Stræben efter \emph{Frihed} som Stræben efter at \emph{herske}). Idet Subjectet under sin Attraa efter Æren tillægger andre Fornuftvæsener et selvstændigt Værd, ligger heri et Element, som, naar det udvikles, ogsaa bringer Mennesket til at foresætte sig \emph{Andres Vel} som sine Handlingers Øiemed.

\anm{Uagtet man ikke med adskillige, især franske, Philosopher kan udlede alle Menneskets Pligter og Dyder af Stræben efter Selvbevaring, saaledes som den her er forklaret, er dog Selvbevaring i høiere Betydning en væsentlig Bestemmelse ved det fornuftige Sindelag.}

\slaw{18}
\emph{Velvillie}, ifølge hvilken Mennesket finder Behag i andre Menneskers Vel og ønsker dets Fremme, kommer først til egentlig Virkelighed i Venskabet og de dermed beslægtede \emph{sociale Forhold}, hvori Velvillien er gjensidig. I sin mindst udviklede Skikkelse grundes Venskabet blot paa den Lyst, som den Ene finder i den Andens Omgang. Det mere forstandige Venskab, hvori den gjensidige udvortes Fordeel er Venskabets Grund, er, hvor der ikke finder anden Grund Sted, kun uegentligen at kalde Venskab. Men i et tredie Slags, det fornuftige og sande Venskab, hvori en fuldstændig Harmonie mellem flere frie Villier finder Sted, er saavel det Behagelige som Nytten bevaret. Derved at dette Forhold (sand \emph{Selskabelighed})\footnote{Forf.\ Msc.\ (Selskabelighedens). F.\ C.\ O.}, hvori Selvbevaringen og Bestræbelsen for Andres Vel ere bragte i Forbindelse med hinanden, \opage{151}antages for Menneskets Bestemmelse, gjøres et betydeligt Fremskridt i den moralske Bevidsthed; men naar denne Forbindelse blot betragtes som en Følge af naturlig Tilbøielighed, er den tilfældig og har Partiskhed til Følge. Den sande Forsoning mellem Selvbevaring og Velvillie finder først Sted i \emph{Retten}, hvori Personen gjelder ikke som Ven, men som Person, og Enhver erholder, hvad der tilkommer ham.

\slaw{19}
\emph{Retten}, den sande Regel for de udvortes Goders Fordeling i det menneskelige Selskab, er et Formaal, der er forskjelligt saavel fra det solitariske, som fra det sociale: vi kalde det et \emph{ideelt Formaal}. Bevidstheden om dette er den frembrydende Bevidsthed om det Godes Idee, men om Ideen under en eensidig Skikkelse. Den strenge Ret kan saaledes stride imod Velvilliens Fordringer, at dens Haandhævelse bliver urigtig, at den strenge Ret bliver Uret og bør vige for \emph{Billighed}. --- Den Erkjendelse, der er nødvendig, for at Retten skal kunne bringes til Bevidsthed, kan ogsaa træde frem som et særeget \emph{ideelt Formaal}, der ogsaa kan tages for Menneskets sande Bestemmelse; Mennesket lever da for \emph{Erkjendelsen og Fremstillingen af det Sande og Skjønne}, for Videnskab og Kunst.

\slaw{20}
Erkjendelsen kan ogsaa for sig eensidigen gjøres til Formaal, saaledes som hos dem, der ansee det \emph{contemplative Liv} for Menneskets Bestemmelse. Øiemedet er her \emph{deels} videnskabelig Erkjendelse, der udvikles saavel i abstraherede Begreber som i rene Fornuftbegreber, \emph{deels} æsthetisk Anskuelse af de Ideer, som ligge til Grund i de tilværende Ting og som ved en høiere Virksomhed af Phantasien, der ei er bestemt ved Erfaringen, fremstilles af Kunstneren eller af dem, hvis Phantasie sættes i Bevægelse ved dennes Værker. Naar den eensidige theoretiske Retning bringes i Forbindelse med de ovennævnte practiske Formaal (de solitariske, \opage{152}de sociale og Rettens Formaal) findes det rigtige Øiemed at være \emph{Handling efter Erkjendelsen}.

\slaw{21}
Til at handle efter Erkjendelse hører ikke blot i ethvert Moment at følge en Maxime, hvorved Livets Øiemed fremmes, men tillige at have almindelige Kundskaber om de tilværende Tings Natur, og Dømmekraft til efter Maximen at anvende Kundskaberne i forekommende Tilfælde. Desuden udkræves en fast Characteer, for ikke ved Hensyn til det Behagelige eller Ubehagelige at lade sig henrive til at handle imod Erkjendelsen. Men Øiemedet, som er Hovedsagen, bliver kun rigtig derved, at den indre Harmonie finder Sted imellem Menneskets lavere og høiere Tilbøieligheder. Deraf udspringer da ogsaa Overeensstemmelse med den udvortes Orden eller den Villie at rette sig efter Objecternes Natur og Forhold til hinanden. Handling efter Erkjendelse udvikler sig saaledes til en \emph{Stræben efter Harmonie mellem alle Menneskets Formaal} og, som en Følge deraf, mellem alle dets udvortes Forhold. Den, der stræber efter, paa en harmonisk Maade at realisere alle disse Formaal, hvortil ogsaa Erkjendelsen selv maa regnes, stræber efter det Gode eller den \emph{fuldkomne Existents}.

\slaw{22}
\emph{Fuldkommenhed} eller fuldkommen Existents er en ved Fornuften bestemt harmonisk Udvikling af alle Menneskets Anlæg og Forhold. At betragte Fuldkommenhed som et Formaal, der for Mennesket er absolut uopnaaeligt, hvoraf Følgen bliver, at ikke den selv, men den i det Uendelige fortsatte Tilnærmelse til den, maa antages for Menneskets Bestemmelse, er en Vildfarelse. I ethvert Moment, hvori Fornuften bestemmer Menneskets Handlemaade, finder en forbigaaende Tilsyneladelse af Fuldkommenheden Sted, og Begrebet om den uendelige Progres har kun forsaavidt \opage{153}Realitet, som Rækken af de Dødeliges fuldkomne Livsyttringer her i Livet ikke kan blive fuldkommen continuerende.

\mchap{Tredie Capitel.}{Om det Gode.}

\begin{center}\textbf{A.\ Om Dyden.}\end{center}

\slaw{23}
Det \emph{fuldkomne Liv}, som i det Foregaaende blev erkjendt for Menneskets Bestemmelse og er det Samme som det Gode, bestaaer i \emph{Fornuftens practiske Virksomhed}, for hvilken \emph{de naturlige Formaal} ere Midler; men da det Gode kun ved dem kan komme til Virkelighed, ere de selv \emph{relative Goder}. Til disse høre saavel de \emph{legemlige Goder} (Sundhed, Styrke, Smidighed) som den \emph{udvortes Lykkes Goder} (Eiendom, Ære \&c.); disse udgjøre tilligemed den indvortes Fuldkommenhed, der ei kan finde Sted, hvor disse Goder aldeles mangle, det \emph{fuldstændige Gode}. Men det \emph{høieste Gode} er den indvortes Fuldkommenhed eller den egentlige Dyd.

\slaw{24}
\emph{Dyden} er den indre Tilstand, i hvilken alle Menneskets Tilbøieligheder ere bragte under Fornuftens Herredømme. Jo mere Dyden er sand Dyd, desto mere har Fornuftens høiere Kraft saaledes gjennemtrængt den naturlige Tilbøielighed, at \emph{Dyden er bleven Mennesket til anden Natur}, eller at de fornuftige Handlinger stemme med dets Tilbøielighed. Men denne Sindelagets Fuldkommenhed er en ved Øvelse erhvervet Færdighed, og ikke nogen blot naturlig Tilstand.

\slaw{25}
Den fornuftige Erkjendelse, som er en Side af Dyden, kan vel fremmes ved Andres Underviisning; men \opage{154}beroer dog tilsidst --- som al Erkjendelse --- paa Menneskets frie Selvvirksomhed, og det saameget mere, som det Gode ikke kan klart erkjendes uden af den, der selv frivilligen har udøvet det. --- Den Øvelse i Dyden, som udkræves, for at den kan blive til Færdighed, kan ogsaa fremmes ved Optugtelse, forsaavidt som Mennesker af Andre kunne tvinges til at indrette deres udvortes Handlinger efter Fornuftens Lov. Men denne udvortes Handling, overeensstemmende med Fornuften, kan kun være et Befordringsmiddel for Dydens Udvikling. Dyden selv som indre Fuldkommenhed kan kun erhverves ved Subjectets frie Selvvirksomhed.

\slaw{26}
Dyden er ingen blot \emph{hvilende Evne}, men en \emph{virksom Kraft}, der lægger sig for Dagen i udvortes Handlinger. Derfor kan ingen kjende sin egen eller Andres Dyd uden af de Handlinger, hvori den træder frem. Al anden formeentlig Bevidsthed om egen Dyd er deels \emph{Selvbedrag}, deels \emph{moralsk Phantasteri}.

\slaw{27}
Forsaavidt som Dyden ei aldeles er modsat Menneskets naturlige Tilstand, men er dennes Bestemmelse ved Fornuften, viser den sig i sin udvortes Fremtræden som en \emph{Middelvei mellem to Extremer}, idet enhver naturlig Tilbøielighed og dens udvortes Yttring i Sammenligning med Fornuftens Fordringer kan være for stærk eller for svag. Men det er kun i sin udvortes Tilsyneladelse, i Tilbøieligheder og Handlinger, at Dyden kan betragtes som indskrænket ved Maal. I sit Væsen er den een Kraft, som ikke kan være for stor.

\slaw{28}
Den i forrige Paragraph omtalte Middelvei mellem to Extremer er ikke eens for Alle. Især kommer herved den forskjellige \emph{Individualitet} i Betragtning, som gjør een Stræ\opage{155}ben mere passende for En, en anden for en Anden. Derfor have Nogle ogsaa med Grund villet inddele Dyderne efter de forskjellige Skikkelser, som Dyden efter Anlægenes Forskjellighed maa antage hos de forskjellige Individer. Saaledes kan man overhovedet adskille de \emph{contemplative} og de \emph{practiske Characterer}, og da disse ere blot adskilte ved den forskjellige Hovedretning, lader den samme Inddeling sig fortsætte paa begge Sider.

\slaw{29}
\emph{Miskjendelse af Individualitetens Natur}, eller planmæssig Afvigelse fra dens Veiledning, frembringer \emph{Usandhed} i Menneskets Liv, hvorved dette mere eller mindre forvandles til et Skinliv. Her er \emph{Affectationen} at mærke, som finder Sted, naar Mennesket, for at opnaae et af de naturlige Formaal, fremviser et andet Forhold mellem sin Person og Tingene, end det, der virkelig finder Sted. Det bringes derved til at yttre paatagne Følelser, Forestillinger og Tilbøieligheder, som det indbilder sig selv at have. Med Hensyn til den større eller mindre Sammenhæng, som Affectationen har med Characteren, kan den snart vise sig som \emph{momentan}, snart som \emph{fast} (idet et skrømtet Træk er optaget i Characteren), snart som en \emph{Vane til Affectation}, der efter Omstændighederne skifter Udseende.

\slaw{30}
Den Dydige besidder den egentlige \emph{Frihed}, da Erkjendelsen, efter hvilken han handler, ikke blot er en Kundskab om Omstændighederne, hvilken ogsaa kan finde Sted hos den, der følger sine Drifter; men en \emph{fornuftig Erkjendelse}, der bestemmer hans Handlingers Øiemed. En saadan Erkjendelse staaer det i \emph{Menneskets egen Magt} at udvikle hos sig, og den udvikles nødvendigen hos den, der ikke ved Vanen, at følge det Behagelige i Stedet for det erkjendte Gode, qvæler dens frembrydende Spirer hos sig.

\anm{Om moralsk Svaghed og Slethed.}

\slaw{31}
\begin{center}\emph{Dyden i sin udvortes Tilsyneladelse.}\end{center}
\opage{156}Vi betragte først Dyden i dens Forhold til de solitariske Drifter. Da Mennesket af Naturen skyer det sandselig Ubehagelige og stræber efter det sandselig Behagelige, viser Dyden sig her deels som \emph{Mod}, deels som \emph{Maadelighed}. Modet er Dyden som Beredvillighed til at gaae Ulyst og Farer imøde, naar Fornuften kræver det. Den ligger mellem \emph{Feighed} paa den ene Side og \emph{Forvovenhed} paa den anden Side. Maadelighed er det fornuftige Maal i Stræben efter sandselig Lyst. Den Last, der er den modsat paa den ene Side som ufornuftigt Overmaal, er \emph{Umaadelighed}; paa den anden Side gives et Extrem i den Skik, at berøve sig selv Livets sandselige Nydelser af misforstaaet Asketik eller Gjerrighed.

\slaw{32}
Saavel i Forhold til den naturlige Sky for det Ubehagelige, som til Attraa efter det Behagelige, viser deels den fornuftige \emph{Vedholdenhed} sig, som bestaaer deri, at Mennesket, uden at afdrages ved hine to Hindringer (Lyst og Ulyst), forfølger en planmæssig fornuftig Virksomhed, deels den fornuftige \emph{Driftighed}, som lægges for Dagen derved, at Menneskets Handlinger følge i Forhold til dets Naturanlæg hurtigt paa hinanden. Forbindelse af Vedholdenhed og Driftighed giver \emph{Arbeidsomhed} eller den planmæssige Driftighed, en Dyd, der ligger mellem \emph{Dovenskab} og \emph{Selvplageri}.

\slaw{33}
I sit Forhold til den naturlige Tragten efter Eiendom viser Dyden sig som \emph{Huusholderiskhed}, der er det fornuftige Maadehold i Erhvervelyst og Udgifter, og har sit Udspring af rigtige Begreber om Eiendom. Den ligger imellem \emph{Gjerrighed} (Vindesyge og Karrighed) paa den ene Side, som kommer af, at man tillægger Eiendom absolut Værdi, og \emph{Ødselhed} paa den anden \opage{157}Side, ifølge hvilken man til mindre vigtige Øiemed gjør saa store Udgifter, at man savner det Fornødne til vigtige Øiemeds Opnaaelse.

\slaw{34}
I Æresdriften ligger Bevidstheden om eget Værd, eller \emph{Selvfølelse}, til Grund for Menneskets Tragten efter at Andre skulle erkjende det. Den fornuftige Grad af Selvfølelse, som meget godt kan bestaae med Ydmyghed, ligger mellem \emph{Stolthed}, under hvilken Forestillingen om eget Værd fremfor Andre formeget fremhersker, og \emph{Pusilanimiteten}, som bringer Mennesket til ufornuftig at begrændse sin Stræben og sine Fordringer. Det fornuftige Sindelag med Hensyn til Æren kaldes \emph{Ærekjærhed} og er lige langt fjernet fra \emph{Ærgjerrighed} og Mangel paa Æresfølelse.

\slaw{35}
Af Ærekjærhed udspringer \emph{Frihedskjærlighed}, ifølge hvilken Mennesket stræber efter en Tilstand, hvori det, uden at tvinges af andre Mennesker, selv kan bestemme sin Virksomhed; en Dyd, der ligger mellem \emph{Tøilesløshed} og \emph{Slavesind}. En mere positiv Form har den af samme Kilde udspringende \emph{Lyst til Indflydelse} eller \emph{Herskelysten}, der fornuftigviis er saaledes indskrænket, at den ikke stiger til Herskesyge, men dog er saa kraftig, at den ikke synker ned til Ligegyldighed for Selskabets fælleds Anliggender.

\slaw{36}
Med Hensyn til de Følelser, der opstaae ved Forestillingen om Andres Tilstand, viser Dyden sig som det \emph{deeltagende Sind}, der ligger mellem \emph{Hjerteløshed} og \emph{menneskekjærligt Føleri}. Med Hensyn til Trangen til at finde Deeltagelse hos Andre, viser det fornuftige Sindelag sig som \emph{Aabenhed}, der paa den ene Side er modsat \emph{Aabenmundethed}, paa den anden Side den overdrevne \emph{Tilbageholdenhed}. Under de \opage{158}sociale Forbindelser, der blot have Lysten til Formaal, træder Dyden frem som \emph{Omgjængelighed}, der ligger mellem \emph{Behagesyge} og \emph{Brantenhed}. Men i de Forbindelser, hvor Nytten er Formaalet, holder Dyden som \emph{Tjenstagtighed} den rette Middelvei mellem \emph{Egennytte} og den forvidt drevne Selvopoffrelse for Andres Skyld.

\slaw{37}
Under sin alvorlige Stræben seer den Fornuftige sig ei blot nødsaget til at modarbeide de Slettes Planer, men bliver undertiden uenig med de Bedre om de Midler, der ere at vælge til de fælleds Formaals Fremme; da viser den rigtige Middelvei sig deels som \emph{Mandighed} eller mandigt Sind, idet man gjør sin fornuftige Villie gjeldende uden at hindres af Behagelyst eller af Frygt, deels som \emph{Føielighed}, idet man retter sig efter Andres Villie, hvor Eens Grundsætninger tillade det. De dertil svarende Overdrivelser ere deels \emph{Herskesyge} og \emph{Vredagtighed}\footnote{Forf.\ Msc.\ har: deels Vredagtighed og Egensind. F.\ C.\ O.}, deels svag \emph{Eftergivenhed} og Opoffrelse af sin individuelle Characteer.

\slaw{38 og 39}
1) I Forhold til Retsdriften viser Dyden sig som \emph{Retsindighed}. Den holder den rette Middelvei mellem den abstracte \emph{Retfærdighed} paa den ene Side, som uden Hensyn til de særegne Omstændigheder, der forvandle den strenge Ret til Uret, forsvarer Rettens døde Regel, og den individuelle \emph{Vilkaarlighed} paa den anden Side, som formeget stræber efter at gjøre den subjective gode Villie gjeldende mod bestemt udtalte Vedtægter og Love.

2) Den rette \emph{Kjærlighed til Kunst og Videnskab} ligger mellem Mangelen paa Sands for disses Ideer og det \opage{159}Afguderi med Kunst og Videnskab, der vil aldeles adskille dem fra Livets practiske Formaal, som blot ere deres nødvendige Forudsætninger.

\slaw{40}
Da det er Fornuften, der bestemmer det rette Maal i Tilbøieligheder og Handlinger, er det Gode ikke godt for Driftens Skyld, men fordi det billiges af Fornuften. Ved Bevidsthed om, at Fornuften i Modsætning til Driften er den høiere og væsentligere Side, bringes Mennesket til at ansee den blotte Tilfredsstillelse af naturlige Begjeringer for \emph{værdiløs} (Stoicisme), og blot at sætte Priis paa den \emph{selvstændige Tænkning}, ved hvilken Mennesket føler sig frit og uafhængigt af de lavere Formaal. Efterat dette Standpunct er naaet, kunne de practiske Formaal ei erkjendes for at være Goder, uden at deres Indhold bestemmes af den rene Fornuft i Stedet for af Driften. Denne Fordring er især i den christelige Tid bleven gjort til Moralphilosophien (Pligtlæren).

\begin{center}\textbf{B.\ Om Pligten.}\end{center}

\slaw{41}
Paa Pligtlærens Standpunct erkjendes, at Fornuften alene skal bestemme Indholdet af Formaalet for Menneskets Handlinger, og at disse ei bør foretages for at tilfredsstille de naturlige Begjeringer (Heteronomie), men af \emph{Lydighed mod Fornuftens høiere Lov}, som Mennesket ubetinget skal underkaste sig (kategorisk Imperativ). Denne Bevidsthed er modsat den antike Dydslære, ifølge hvilken den blot naturlige Begjering inden en vis Grændse var god og afgav Stoffet, der reguleres af Fornuften. (Uagtet den, der handler af Ærefrygt for Pligten, underkaster sig et strengt Bud, der har Characteren af en høiere Magt, reddes dog Villiens Frihed ved Bevidstheden om, at Fornuften er Menneskets eget inderste Væsen.) \opage{160}Naar denne Synsmaade conseqvent forfølges, maa Fornuftens Virksomhed erkjendes for det eneste Formaal, der bør attraaes for sin egen Skyld, og de naturlige Formaal maae betragtes som \emph{relative Goder}, der ikkun have Værd som Midler til dette Formaals Fremme.

\anm{Om Inddeling af Pligter. Ubetingede Pligter, ligegyldige Handlinger.}

\slaw{42}
\emph{Pligtcollision} finder Sted, naar Fornuften synes at paabyde flere Handlinger, som staae i et saadant Forhold til hinanden, at den enes Udførelse gjør den andens umulig. Om der nu end for Pligtcollisionen kunne gives flere \emph{Klogskabsregler}, hvorved Collisionernes Antal vil formindskes, kunne dog ingen almeengyldige Regler gives, hvorved nogen for alle Tilfælde gyldig \emph{Rangforordning} mellem de almindelige Classer af Pligter kan bestemmes. Af de tre Hovedclasser kunne hverken Pligterne mod os selv eller mod Andre eller mod Gud sættes som de, der i alle Tilfælde skulle foretrækkes. Ligesaalidt kan en almeengyldig Orden bestemmes for de underordnede Pligter i nogen enkelt af disse Hovedclasser. Forstanden maa da opgive sit Forsøg paa at give Pligtlæren sin Fuldendelse i Detaillet, og der bliver kun den Udvei tilbage, at arbeide paa \emph{Overvindelse af Sandseligheden}, der hindrer Pligtbudenes rigtige Anvendelse i de enkelte Tilfælde.

\slaw{43}
Ved \emph{asketiske Øvelser} søger Mennesket at \emph{overvinde Sandseligheden}, som en Hindring for Fornuftens Virksomhed. \emph{Bodsøvelse} og den i Eensomhed anstillede \emph{Selvprøvelse} anvendes som Midler til at erhverve Dyden; men især anbefaler Asketiken idelig at henvende Tanken i Bøn til Guddommen --- Pligternes Kilde. Ved disse Dydsmidler, især ved det sidste og fuldkomneste, stræber Mennesket at tilveiebringe \opage{161}en saadan Overeensstemmelse mellem sig selv og Pligtloven, at det af sig selv som enkelt Individ kan øse Pligtbudets Kjendelser. Det handler da efter sit eget Hjertes Tilsagn eller efter sin Samvittighed.

\begin{center}\textbf{C.\ Om Samvittigheden.}\end{center}

\slaw{44}
\emph{Samvittigheden} erkjender umiddelbart, hvilke Handlinger der paabydes af Fornuften, og Subjectet er i Kraft af den aldeles vis paa, at det under sin Handling er i Overeensstemmelse med det Gode (det har en god Samvittighed). Den stiller sig endog over Morallovens bestemte Forskrifter, fordi den i det enkelte Tilfælde er sig bevidst at have Ret til at afgjøre deres relative Gyldighed. --- Men Samvittigheden kan dog være vaklende, ja vildfarende, da den jo endogsaa kan gjendrive sig selv, og det maa følgelig komme an paa Handlingens Indhold, om den skal dømmes at være god eller ikke. Saaledes kan da Samvittigheden ved at gjøre sig selv til Lovens Kilde komme til at bestemme sig tvertimod Moralloven eller til det Onde.

\slaw{45}
Det Onde kan udøves saaledes, at Subjectet har Bevidsthed om Pligtens Bydende, men dog handler derimod; da siges det at handle med \emph{ond Samvittighed}, men det kan og med \emph{Hykleri} for sig selv og Andre besmykke sin onde Handling derved, at det lader en eller anden Bevæggrund gjelde som antagelig og god Grund derfor (\emph{Probabilismus}). Idet en saadan antagelig Grund kan gives for enhver tænkelig Handling, bliver der ikke Noget tilbage af særskilt godt Indhold, men kun den almindelige gode Hensigt.

\slaw{46}
Den gode Hensigt antages da at hellige Midlet, og \opage{162}da en saadan Hensigt kan udhæves ved enhver slet Handling, er det formeentlige Gode herved aldeles forvandlet til Ondt. Et Formaal vælges, fordi det forekommer Subjectet godt. Det handler efter sin Mening, og det gjør hos det Intet til Sagen, om det, der efter dets Mening er godt, i sig selv er ondt. Da med denne Beroen ved den subjective Mening det i og for sig Gode er fornegtet, er derfra en let Overgang til den \emph{moralske Ironie}, ifølge hvilken man holder sin individuelle Villie for det Høieste og følgelig troer sig hævet over Moralloven. Hermed har den subjective Frihed opnaaet sin formelle Fuldendthed, men mistet alt objectivt Indhold.

\slaw{47}
Det er i de sidste Paragrapher viist, at Bestræbelsen for at erkjende Menneskets Bestemmelse blot fra Individets Standpunct er frugtesløs og kan i sin Eensidighed føre til fornuftstridige Principer. Det Gode skal ei erkjendes blot ved subjectiv Tænkning, men findes i sin Virkelighed og Udvikling i Staten, og Individet opfylder kun sin Bestemmelse ved at virke som Led i et saadant fornuftigt organiseret Selskab. Derved er den individuelle Frihed ei ophævet eller indskrænket; thi Individet maa kunne gjenfinde den samme Fornuft, som udvikler sig hos det Selv, i Statens Love, Vedtægter og Indretninger. Det finder deri kun en Grændse for sin Vilkaarlighed, idet den fornuftige Friheds Grundformer, hvilke den sande Samvittighed ikke kan byde Mennesket at overskride, deri have deres objective Tilværelse. I Statslivet erholder de naturlige Drifters Tilfredsstillelse sin sædelige Betydning, idet de blive til Midler for Statens Liv og Fremgang, og den frieste Udvikling af de individuelle Anlæg fremmer Selskabets Fuldkommenhed og Lyksalighed.

\ombreak

\end{document}
